\section{Introduction}
\label{sec:intro}

Graph neural networks are the default tool for relational data, and their cost profile is becoming
the binding constraint. Message passing touches every edge at every layer, so inference on a large
graph is memory-bandwidth bound rather than compute bound, and streaming or embedded
deployments---sensor meshes, embedded screening, on-device recommendation---must meet a latency
budget that changes with the device and the moment. Spiking neural networks (SNNs) offer a
structurally different execution model: computation is event-driven, neurons integrate and fire, and
a step costs a sparse set of synaptic events rather than a dense tensor product. Whether that also
saves \emph{energy} depends entirely on how many spikes the mapping emits, which is an empirical
question and not one this paper assumes; the property conversion does guarantee is different, and it
is the one we build on.

There are two ways to obtain such a network. The first is to train it directly: unroll the spiking
dynamics over $T$ timesteps, replace the derivative of the spike with a surrogate, and run
backpropagation through time. This works, and it is what essentially all spiking graph networks do
today. It is also expensive in a way that is easy to underestimate: the network must be retrained
for every latency budget, because $T$ is baked into the unrolled graph, and training cost grows
linearly in $T$ on top of a backward pass that is already $T$ times a forward pass. A deployed model
therefore cannot adapt when the latency budget changes---it must be retrained, off-device, by
someone who has the training pipeline. The second route is \emph{conversion}: train an ordinary ANN,
then replace its rectifiers with spiking neurons and its activations with firing rates. Conversion is
attractive precisely because it decouples training from latency---one training run, deployable at
any $T$, with no backward pass anywhere near the device---and in vision it has matured to the point
where the accuracy gap at modest $T$ is a fraction of a percentage point.

For graphs, conversion \emph{methods} already exist --- a training-free conversion of a text
classification GNN \citep{hwang2024onespike} and a tactile-graph conversion with structural
retraining \citep{zhang2025quantization} --- but we are not aware of a published \emph{analysis} of what the
conversion error is made of on a graph, or of any per-layer window allocation in graph conversion.
That analysis is
the subject of this paper, and the borrowing of vision recipes fails for a reason that
is specific to message passing. In a convolutional network the input to a unit at step $t$ is a
fixed tensor: every channel of every receptive field is available simultaneously. In a graph, the
input to a unit is a \emph{sum over neighbours}, and each neighbour fires at its own time. The unit
therefore integrates a sequence of partial sums whose addends arrive out of order, and the
instantaneous current at step $t$ is the ANN pre-activation with its terms reordered in time. We
call this \textbf{neighbourhood temporal misalignment}. The name is deliberately qualified:
\emph{temporal misalignment} has been used for feed-forward conversion, where random rearrangement
of spikes \emph{across layers} can improve accuracy and can be exploited with stochastic neurons
\citep{bojkovic2025temporal}. Ours is a different object. It lives \emph{inside a single
aggregation}, over a set of neighbours that share no clock, so it is a property of the graph rather
than of the network's depth---and, as we show in \secref{sec:ablate}, the two phenomena lead to
opposite conclusions about whether stochasticity helps. \figref{fig:framework} summarises the
setting: one trained GCN, a latency-agnostic deployment, and the neighbour-level clock skew that the
rest of the paper quantifies.

Those artefacts are worth stating precisely, because they set the shape of the method, and two of
them are integer effects that conversion necessarily introduces. First, a spiking neuron driven by a
constant current emits $\lfloor Tc/\theta \rfloor$ spikes, so its rate represents $c$ with an error
that is \emph{systematically negative} and of size $1/(2T)$; in the vision literature this is
corrected by adding a per-step bias, and we show that on graphs the correction has a closed form.
Second, if the output layer is decoded by \emph{counting} spikes, the decision is quantised to
multiples of $1/T$ and---more damagingly---the logits are clipped at one spike per step, which
destroys the ordering of large logits long before the rate has converged. A third artefact is
inherited from vision but inverts on graphs: conversion normally rescales activations so that they
fit under the threshold, and this rescaling assumes that averaging operators are row-stochastic.
Symmetric GCN normalisation is not: the row sums of $D^{-1/2}(A+I)D^{-1/2}$ are $\rhon$, which we
measure at up to $5.75$ on Cora, so the usual global scale factor is simply the wrong object.

Writing misalignment down makes the rest analytic. \textbf{Conversion error in a spiking graph
network is deterministic}: it is a function of the trained weights and the latency, not a random
variable to be averaged over seeds. The closed form, the spectral-gain law, the order invariance, and
the boundary at which the closed form stops holding all follow from that claim, as does the one
prediction that fails. Our contributions are:

\begin{enumerate}\itemsep2pt
\item \textbf{A conversion formulation for graph node classification}, with an exact error identity,
$\Btwohat - \Btwo = \Ahat(\Eone W_{2})$ where $\Eone = \cntone/T - \mathrm{relu}(\Bone)$, that
separates the conversion error into a term the neighbourhood averages and a term it cannot
(\secref{sec:align}). No approximation is made in deriving it, which is what allows the theory and
the implementation to be checked against each other exactly (\secref{sec:theory}).
\item \textbf{A parameter-free bias.} In the population, $\mathbb{E}[\Eone] = \beta - 1/(2T)$, so
$\betastar = \theta/(2T)$ is the unique per-step bias that cancels the systematic error
(\secref{sec:bias}). The bias is inherited from forward-bias conversion; what is new is that on
graphs it admits a closed form, which removes calibration altogether, and that it beats the
train-mask-selected bias grid (step $0.05$) in all nine graph-by-latency cells we test.
\item \textbf{A membrane readout} that decodes the \emph{accumulated} layer-2 current rather than the
spike count (\secref{sec:readout}). Because $\mathrm{mem} = \cntone\theta + u = \sum_t
\mathrm{inst}_t$ holds exactly, the readout costs nothing in spikes, yet it removes both the $1/T$
decision grid and the unit-clipping ceiling: at $T=8$ it is worth $3.7$ points on Cora and $6.2$ on
CiteSeer, largest exactly where the clipping binds hardest. The same identity has a structural
consequence we did not expect and now test directly: the accumulator contains no $t$, so the
membrane decoder is exactly invariant to the \emph{order} of upstream spikes while the count
decoder is not (\secref{sec:order}). That turns a preference into a mechanism, and it is what
separates our error from the cross-layer temporal misalignment of \citet{bojkovic2025temporal}.
\item \textbf{A margin-recovery theorem} (\secref{sec:thm}) giving the number of timesteps needed to
preserve the ANN's decision, $T = O(\anorm \lVert DW\rVert_{2} / \Deltan)$, verified cell by cell
with a theory-versus-simulation gap of exactly zero, together with a quantitative account of where
the bound is loose (\secref{sec:tight}).
\item \textbf{A depth law, and a negative result.} The identity extends to $L$ layers
(\eqnref{eq:recursion}), where it says the output error is amplified by the cumulative spectral gain
$G=\prod_l\lVert M_l\rVert_{2}$ of the trained ANN. Measured over nine cells the amplification follows
$G^{0.95}$ with $r=0.990$, so the latency a deep model will need is computable from its weights
before any simulation (\secref{sec:depth}). Two natural repairs---gauge-normalising each layer, or
rescaling the readout---do not work, and we say why: the spike threshold is an absolute scale, so
$G$ cannot be reparametrised away. Conversion therefore favours shallow, sufficient-width models,
and width is the lever that does work: the layer-1 error norm grows as $\sqrt{H}$ while the logit
margin grows with it, so the fidelity gap stays near zero up to $H=64$ even though the ANN gains
$0.8$--$1.9$ points.
\item \textbf{A conversion that needs no simulation.} For a layer driven by a constant current the
count over $T$ steps is $\lfloor T(B+\beta)/\theta\rfloor$ in closed form, so the two-layer
conversion can be scored with one sparse matmul instead of $T$ steps of spiking. Measured over $105$
cells the closed form and the simulator agree to $0.08$ points on average and $0.40$ at worst, which
is float32 drift on one or two test nodes (\secref{sec:closed}). At $L=3$ the same closed form fails
by $3.5$ points, and the reason is the finding rather than an implementation detail: the reset makes
an intermediate layer's \emph{count} depend on the order of its input even though the output
membrane does not.
\item \textbf{What the laws do and do not predict.} Fitting the single constant of
\eqnref{eq:sigma} on some cells and scoring others, the machinery ranks candidate models almost
perfectly ($\rho=0.967$ under leave-one-configuration-out) but misses the absolute gap by $4.6$
points on average, and the ablation the mechanism implies --- scrambling the per-node error scales
--- does not hurt. We report this as a failed prediction and claim only the ranking
(\secref{sec:ranking}).
\item \textbf{A four-arm Pareto under one protocol.} Conversion, a directly trained spiking GCN, the
rate-coded published formulation, and the published-architecture formulation, all at seven seeds on
five latencies. No arm dominates the conversion anywhere; the cost ordering depends on the
accounting convention --- under a static synaptic-update proxy the conversion is the cheapest arm at
matched accuracy on all three graphs, while under event-driven counting a linear-readout external
arm is cheaper at low latency on two of them --- and the published-architecture formulation stays
more accurate on two graphs (\secref{sec:pareto}).
\item \textbf{A constructive latency decision in the deep regime}: where uniform budgets explode
with the depth law, a validation-only per-layer window search (the idea family is established prior
art; the graph instantiation and regime measurement are ours) keeps the accuracy band with
$22.7$--$93.4\%$ fewer layer-updates on six of eight held-out four-layer cells --- with the
pre-registered mixed-depth screen disclosed as not met --- and three weight-only routes are
reported as measured transfer failures (\secref{sec:allocation}, \secref{sec:transfer}).
\item \textbf{An empirical study at seven seeds} across six graphs, three aggregation operators, three
hidden widths and two trained spiking baselines (\secref{sec:exp}). Converted networks reach the ANN
to within $1.5$ points at $T \le 8$ on all three citation graphs; they never need more timesteps than
a tuned directly-trained spiking GCN to reach a given accuracy, are two-fold earlier on Cora, and are
the only arm that reaches the band at all on CiteSeer; they are clearly ahead of the rate-coded
formulation used by published spiking-GCN work. Every one of those comparisons is reported against a
measured run-to-run floor rather than against a star (\secref{sec:baselines})---and the method
degrades in the \emph{predicted} way when the aggregation
operator, the threshold, the width or the decoder is changed.
\end{enumerate}

We state the boundaries of these claims early, because they are load-bearing. The scale experiment
(ogbn-arxiv, $169{,}343$ nodes) is a \emph{fidelity} test: we ask whether conversion preserves a
given ANN, not whether our ANN is the best available on that dataset. Our efficiency evidence is
spike counts and timesteps, not joules measured on neuromorphic hardware. And our strongest baseline
is a faithful re-implementation of a published spiking-GNN \emph{formulation}, not a reproduction of
a published codebase. Each of these is revisited in \secref{sec:boundaries}.
